\documentclass[11pt]{article}
\usepackage{amsmath,amssymb,amsthm,geometry}\geometry{margin=1in}
\newtheorem{definition}{Definition}\newtheorem{theorem}[definition]{Theorem}\newtheorem{proposition}[definition]{Proposition}
\title{Flat scope-library descent and exact source boundaries}
\date{16 November 2026}
\begin{document}\maketitle
\paragraph{Audit.} Native flat-library promotion theorem. It constructs an actual operational/schema descent, not merely a quotient of observed bodies. Its domain is the common-boundary actual-limit branch, all admitted joint laws, stable occurrence addresses, deterministic port-free seeds, and separate full terminal polynomial grades. It does not cover arbitrary boundary graphs, tensor/limit mixtures, random silent initializers or a strong inclusion of the original source. Neither earlier reviewed source is edited.

\paragraph{Exact dependencies.} The global theorem used is \emph{Free Global Scope Plan Algebras v2}, SHA256 \texttt{5568c1ae5781688c1dad959e6b8dd97d6685dfb43a660139d3b9a9dfc6e61390}, independently reviewed at that declared conditional scope by \texttt{53a07148e608bb700d26c77453c98818e0a00bfa5994e6d08783fc7102faa8f1}. The corrected standalone operational statement is \emph{Boundary Row Algebras Command Indexed v2 Corrected}, SHA256 \texttt{7c2b6f46bbcb7923b2451921c0932daa47d3d0aa94987aae49e82b0fb91cbfbb}; it incorporates the original review's explicit interface pins with the explicit interface pins stated there. The source assembly is SHA256 \texttt{bbc0006ed160be930687306e2f57b6ab57d111e93b13e6c92b4f80dc881d60cd}; the extension contract is \texttt{642e10beb1d4d0d16703e3d4023c2d69c26631a989d8be3505a03be0d226a0a4}; and the earlier source/feedback contract ledger is \texttt{ab12a98ab64b36f1cf9311a710425016a14795e6446dc1bf55e4ad3efa5fb902}. These were read, not inferred from filenames.

\section{What the source actually specifies}
The assembly's section on scope constructors states that a finite schema includes primitive command schema labels, the boundary diagram and the coefficient mode; repeated occurrences have distinct formal positions. It retains every formal schema and unused shape, and library maps are induced by primitive-schema and actual boundary maps. Its rebracketing proof uses the same flat primitive diagram on both sides. This is not a doctrine identifying arbitrary collections merely because they are isomorphic, or identifying arrows because realized bodies agree.

The text does not say that these listed fields are the \emph{only} permissible schema metadata. In particular it does not require identification of every passive root frame. An implementation retaining complete scope-plan provenance is therefore an enrichment compatible with the specified fields, but its extra distinctions must be declared. A minimal flat-diagram implementation retains only the declared ordered addressed active primitive diagram, its typed boundary and coefficient mode. That implementation does identify distinct plans whose passive-frame metadata differs while the active diagram agrees. The theorem below concerns exactly this independently specified implementation. No mandatory quotient is attributed to the source beyond its explicit definition.

For example, take one active root command schema $s$ enabled both in a root state with no recursive ports and in a root state with one recursive port and a wholly passive child. With the same root occurrence address, world boundary and limit mode, the two flat active diagrams each have the single root vertex labeled $s$. The free schemas differ because their frames/child slots differ. This is a genuine alias in the minimal flat implementation. It is not an equality of configurations or menus: those remain distinct, as do their future behaviors.

\section{The independently fixed flattening map}
Use the sorts $(w,u,e)$ of the global theorem. Strengthen each formal root frame to carry the actual stable identifiers of its recursive ports, as well as their input colors and planar order. Denote these keys by $a_1,\ldots,a_n$. They are not their current ranks: adding a passive sibling must not rename an old occurrence. A primitive address is a finite sequence of such keys. The root address is the empty sequence. The planar order of active addresses is retained separately; no lexicographic order on a fixed key alphabet is assumed to reproduce a changing port order.

Define $E$ on every formal free plan, realized or not. Its output consists of the sort $(w,u,e)$, the ordered finite set of active stable addresses, the primitive schema at each address with its input type, the complete schema-level full-family boundary incidences, and the declared coefficient mode, here limit mode. Include any further state-independent coefficient-diagram annotations actually required by the selected source library. The fixed boundary $e$ includes its entire raw carrier, actual-world-arrow projection, context boundary coefficients, full writer boundary and body. Neither $q$ nor its boundary marginal $\nu$ determines this index.

Recursively, an active root contributes $(\varnothing,s)$; a passive root contributes no active vertex. For an active child plan at port key $a_j$, prefix every active address by $a_j$. Concatenate these blocks in the root frame's planar port order, retaining the root-first convention. A passive child contributes no active vertex, but operationally still receives its same boundary-driven rebase. Define $\overline A$ as the image of $E$ on all formal terms, not the image of actual commands or positive-support executions. Every image diagram has at least one active vertex.

\begin{proposition}[Substitution congruence]
The relation $t\sim t'$ iff $E(t)=E(t')$ is a many-sorted congruence for the formal root constructor. It retains stable active addresses and all declared active diagram annotations.
\end{proposition}
\begin{proof}
Fix a root frame and its root active/passive choice. Replacing a child term by one with the same $E$ leaves its ordered address list, typed labels and boundary incidences unchanged. Prefixing by the same stable port key therefore gives the same parent list and diagram. Distinct child blocks have disjoint prefixes. Thus replacing any or all children by equivalent terms preserves the parent diagram. This argument uses no realizations, probabilities or successor states and applies to all unused formal terms. Entirely passive children are the fixed boundary marker, not an additional term needing a representative. Frames with different passive profiles can have the same $E$, but inserting either as a child still prefixes the same active diagram by the parent's fixed key. Flattened ranks would not justify this argument and are not used.
\end{proof}

For $D\in\overline A$, define $\overline Q_D$ to be the actual full-family limit of its retained active diagram over $H_e$. For every plan $t$, define
\[
 \theta_t:Q_t\xrightarrow{\sim}\overline Q_{E(t)}
\]
by its full primitive projections at their stable addresses. Every passive factor is the terminal object of the slice, hence contributes only the already retained boundary projection. The iterated-limit and flat-limit universal properties give mutually inverse arrows. These are actual arrows on all shapes and their coefficient objects, including those never realized by a body. The empty active set is excluded at the whole-command level; neutral passive subblocks still use $H_e$, not an ambient tensor unit.

All diagrams comparing $\theta$ with child substitution commute: every route has the same primitive-address projections and boundary projection. In particular $\theta_{t'}^{-1}\theta_t$ for $t\sim t'$ satisfies identity/cocycle laws and is compatible with all root constructors. This proves full-library coherence without invoking equality only on actual body images.

\section{An operational quotient with no raw or law quotient}
Let $F^{\mathrm{plan}}$ be the reviewed finite-feedback carrier with its entire free plan library. Define $F^{\mathrm{flat}}$ to have exactly its configurations, seeds, passive action, external ports, completion markers, command set, full raw carrier for each command, supplied law $q$, actual common update, post-state, and actual context coefficient. Its terminal collection remains exactly the same entire colored polynomial tree collection. Change only the writer label of a command from $t$ to $E(t)$ and its full writer family to $\overline Q_{E(t)}$, transporting the body by $\theta_t$.

There is a native map $q_F:F^{\mathrm{plan}}\to F^{\mathrm{flat}}$ whose state, menu, raw and context maps are identities and whose whole terminal grade map is identity. Its global schema map is $E$ and its full-family arrows are $\theta_t$. It does not identify commands with different joint laws, merge passive states, remove raw zero atoms, change refusal/completion, or merge occurrence addresses. A shared schema name across two commands is not the same as identifying those commands.

The constructor $\overline\iota:T_PF^{\mathrm{flat}}\to F^{\mathrm{flat}}$ grafts states and scopes as before. Its global schema function sends a formal root frame with flat child diagrams to their prefix-and-union diagram. The congruence proposition proves this independent of any representative child plan. Its writer arrow is the actual limit-of-limits comparison from the product of the root and full child families to the complete flat active family. At commands its law is the full retained outer joint flattened by the raw coordinate bijection, not the product of child marginals. The actual context arrow is the same flat coefficient comparison as before. Active post-states, passive rebases and birth seeds are unchanged. Thus this is a native algebra, and
\[
 q_F\,\iota=\overline\iota\,T_P(q_F)
\]
holds on every field. On full writer families it follows from the identical primitive projections on all formal schemas; on raw laws it follows from retaining the same entire joint under the same tuple bijection.

\section{The exact equational algebra category}
Fix the syntactically specified equations $t=t'$ whenever $E(t)=E(t')$, and equip them with the coherent full-family comparison $\theta_{t'}^{-1}\theta_t$. These equations are determined by independently typed finite active diagrams before choosing any target algebra. Define a \emph{flat-library algebra} to be a native $T_P$-algebra whose global schema algebra satisfies these equations, including their full-family arrows. Equivalently, the interpretations of the two sides have equal schema names and their full-family arrows agree after this specified comparison. This is a presentation by schema equations, not an assumption that any closed operational feedback comparison already commutes. It imposes no new condition on state functions, joint laws or raw transitions; their equations remain to be proved by the native fold construction.

\begin{theorem}[Initial flat-library feedback]
$F^{\mathrm{flat}}$ with $\overline\iota$ is initial among flat-library native algebras. The corresponding fold preserves every configuration, command, complete raw row/law, actual coefficient, full writer family and separate terminal grade at the previously declared native comparison strength.
\end{theorem}
\begin{proof}
The free-plan initiality theorem gives a unique full native fold into any target algebra. Its global schema function and whole-family arrows interpret the formal schema terms. For a flat-library algebra the defining diagram equations therefore make these maps constant on $E$-fibers, with the specified coherent family comparisons. They define one map on every $\overline Q_D$, independent of a representing term. Configurations and commands were not quotiented; use exactly the free fold's state, menu, raw, context and terminal fields. Its body square follows by composing the descended family arrow with $\theta_t$. Hence it is a native fold from $F^{\mathrm{flat}}$ and obeys the algebra equation by the already established constructor compatibility. Conversely composing any such fold with $q_F$ gives the unique free fold. Surjectivity of $E$ on all formal diagrams, and the isomorphisms $\theta_t$ on their entire families, force equality of global schema and full-family fields, including unused diagrams. All other fields are forced because $q_F$ is identity on them.
\end{proof}

For unrestricted target algebras this statement is false whenever $E$ is genuinely noninjective. The free-plan algebra itself is a permitted target, and its unique fold from $F^{\mathrm{plan}}$ is the identity on all schema names. That identity cannot factor through two distinct terms with equal $E$. This is an obstruction at the global library-map field despite identical configurations, raw laws and coefficient objects. Restricting to the independently specified equation-respecting category is therefore substantive, not redundant. The population algebra from the global theorem belongs to this class: its schema maps are constant at each sort and its full families are units. It supplies an actual correlated, hierarchy-forgetting semantic algebra after descent.

\section{What this does and does not promote under the contract}
This is an E1 mapping property for an independently specified non-rigid native operational category and its flat-library equational subclass. It interprets feedback using actual coupled scopes, passive rebasing and opaque target algebras; it is not solely syntax. The quotient construction shows precisely how the free provenance library relates to the minimal flat source implementation. It supplies E2 unfolding and a native universal fold for deterministic port-free recursive invocation, coupled with E3 whole-scope root decomposition. It is not a characterization of the entire mixed context/content extension, and no E1--E6 completion is claimed.

The original source has independent context/content presentations and an earlier strong initial-row inclusion. The present port-free seed interface is narrower. Moving a source initial row into a paid boot command produces a port-free operational interpretation, but changes the paid scope interface. The source ledger's one-boot versus two-boot enabled-count obstruction prevents promoting that paid delay to an unrestricted horizontal strong or lax source inclusion merely from these feedback results. Thus neither E5 conservativity nor old strong source recovery follows from this new initiality.

At trivial world with identity passive actions, a persistent commuting-event primitive can be represented by local commands consisting of its nonempty event subsets and their full supplied couplings. Partitioning a global selected event set by invocation gives the same raw tuple carrier, law and post-state as the new finite feedback carrier. This is the source's already proved operational specialization. The additional coefficient mode must still be declared: old event tensors are not generally the new slice products. Over an augmentation unit there is the natural directed context map from the tensor to the product, obtained from the component augmentations; it is generally not invertible. Where the monoidal structure is Cartesian and the boundary is terminal, or where all relevant coefficient factors are units, these actual coefficient comparisons are isomorphisms. Scope writer families require the same compatible mode specialization independently; the terminal polynomial grade remains its original tensor construction. Consequently a full native old-feedback isomorphism is warranted only on that coincident-mode restriction, not throughout the original augmented-vector doctrine. In the tensor-to-limit branch the interpretation is directed, with the exact kernel retained.

A further finite-initialization theorem would have to replace selected port-free seeds by full initial raw rows and actual coefficient/body arrows. Admission must be all-raw finite/well-founded at every world and recursive input reachable by the typing, not just at the original seed world. Initialization witnesses and their joint law with the triggering scope must be retained; selecting fresh independent draws is not automatic. Root decomposition must preserve those witnesses, world timing and coefficients, and algebra folds must recurse the initial-row arrow as well as the paid-row arrow. No equation rebasing an old seed into a fresh seed is available. These are concrete residual inputs and proof obligations, not implicit corollaries of the present deterministic-seed fold.

Retained-state instantaneous cyclic closure, seeded delayed recurrence and more general iteration remain distinct. This theorem supplies the stated recursive-invocation operation and its actual interaction assembly; it does not supply an instantaneous trace or compare all retained-state closures. Infinite histories and accumulated infinite content are likewise not needed to prove this finite-state-carrier initiality, but their broader contract obligations are not removed.
\end{document}
